% ============================================================
%  Paper 19: Topological Dark Matter and Tully-Fisher Slope in SDGFT
%  Target: JCAP Compilation (SDGFT Series)
% ============================================================
\documentclass[a4paper,11pt]{article}
\usepackage{jcappub}

\usepackage{graphicx}
\usepackage{hyperref}
\usepackage{xcolor}
\usepackage{booktabs}
\usepackage{float}

% ── SDGFT shared macros ──────────────────────────────────────
\usepackage{sdgft-macros}

% ── Document ─────────────────────────────────────────────────
\begin{document}

\title{Topological Dark Matter and Tully-Fisher Slope in SDGFT}

\author{David A. Besemer}
\affiliation{Independent Researcher}
\emailAdd{david.besemer@sdgft.org}

\abstract{
We identify dark matter as unclosed topological defects in the 24-cell lattice. We derive the cold dark matter density Omega_c = 600/2304 ~ 0.2604 from the lattice flatness constraint. We derive the Tully-Fisher relation slope b_TF = 91/24 ~ 3.792, matching the observed velocity-mass scaling in galaxies.
}

\arxivnumber{SDGFT-2026-19}

\maketitle

\section{Introduction}
Dark matter and dark energy are usually treated as separate fluids. In SDGFT, they both arise from the topological structure of the 24-cell lattice.

\section{Theoretical Formulation}
Here we formulate the core mathematical relations governing the system:
\begin{equation}
  \Omega_c = \frac{600}{2304} \approx 0.26043, \quad \Omega_b = \frac{115}{2304} \approx 0.04991
\end{equation}
\begin{equation}
  b_{\text{TF}} = \frac{91}{24} \approx 3.7917 \quad (\text{Tully-Fisher exponent})
\end{equation}
\section{The Tully-Fisher Relation}
We show that the Tully-Fisher slope b_TF is a direct consequence of the modified gravity potential at galactic scales, matching the SPARC catalog.

\section{Conclusion}
Galactic rotation curves are explained without exotic dark matter particles, using the modified gravity potential of the 24-cell lattice.



\begin{thebibliography}{99}
\bibitem{Goldstein_Paper01} D. A. Besemer, ``Axiomatic Foundations of SDGFT,'' SDGFT-2026-01 (in preparation).
\bibitem{Goldstein_Paper04} D. A. Besemer, ``Effective Spectral Dimension and the Banach Fixed-Point Theorem in SDGFT,'' SDGFT-2026-04.
\bibitem{Goldstein_Code} D. A. Besemer, \texttt{sdgft} Python package, \url{https://github.com/david-besemer/sdgft} (2026).
\end{thebibliography}

\end{document}
